
Il
existe une extension finie @@M0@@ de @@M1@@ telle que @@M2@@ ait toutes ses extensions résiduelles triviales. Si @@M3@@ est un schéma intègre fini et plat sur un ouvert @@M4@@ de @@M5@@, de corps de fonctions @@M6@@, il existe un ouvert @@M7@@ de @@M8@@ tel que
@@M9@@
au-dessus de @@M10@@~; quitte à restreindre @@M11@@ et @@M12@@, on peut supposer @@M13@@ et @@M14@@. Comme l'assertion à démontrer est locale pour f.p.q.c., on peut remplacer @@M15@@ par @@M16@@. Enfin on applique 3.1.

\textbf{3.6. Remarque.} Il résulte de même, compte tenu des suites exactes @@M17@@, que pour tout @@M18@@
@@M19@@
est un groupe unipotent.

\section*{4. Démonstration de II* : la partie la plus délicate de la démonstration}

Rappelons l'assertion II*~: le théorème 1.1 lorsque @@M20@@ et @@M21@@ sont intègres sur @@M22@@ et la suite (*) @@M23@@ est exacte~; en appliquant le théorème de platitude générique EGA IV 6.9.1 et EGA III 7.8.8, on se ramène à prouver la

\textbf{Proposition 4.1.} \uline{Soient} @@M24@@ \uline{un schéma noethérien intègre,} @@M25@@ \uline{et} @@M26@@ \uline{des morphismes propres de présentation finie,} @@M27@@ \uline{un} @@M28@@\uline{-morphisme surjectif et} @@M29@@ \uline{le morphisme canonique. On suppose que}
\begin{enumerate}[label=\arabic*)]
\item @@M30@@ \uline{sont plats sur} @@M31@@.
\item @@M32@@, @@M33@@ \uline{universellement.}
